%% kapa: exam=1 round=34 year=2023 period=1 form=A subject=real_estate subject_ko=부동산학원론 seq=26 number=106
%% source: Q-Net 감정평가사 기출문제 (cst00309) artlSeq=5211789 file=raw/1008/5211789/2023년도 제34회 감정평가사 1차시험 1교시.pdf
%% conversion: kordoc 4.24.51 -> tools/gen1.py (auto); review=proofread-sonnet; flags=none
\begin{question}{106}
부동산개발사업의 위험에 관한 설명이다. (\quad)에 들어갈 내용으로 옳은 것은?
\begin{framedtext}
○ ( ㄱ )은 추정된 토지비, 건축비, 설계비 등 개발비용의 범위 내에서 개발이 이루어져야 하는데, 인플레이션 및 예상치 못한 개발기간의 장기화 등으로 발생할 수 있다. \par
○ ( ㄴ )은 용도지역제와 같은 토지이용규제의 변화와 관계기관 인허가 승인의 불확실성 등으로 야기될 수 있다. \par
○ ( ㄷ )은 개발기간 중 이자율의 변화, 시장침체에 따른 공실의 장기화 등이 원인일 수 있다. \par
\end{framedtext}
\begin{choices}
  \item ㄱ: 시장위험, \quad ㄴ: 계획위험, \quad ㄷ: 비용위험
  \item ㄱ: 시장위험, \quad ㄴ: 법률위험, \quad ㄷ: 비용위험
  \item ㄱ: 비용위험, \quad ㄴ: 계획위험, \quad ㄷ: 시장위험
  \item ㄱ: 비용위험, \quad ㄴ: 법률위험, \quad ㄷ: 시장위험
  \item ㄱ: 비용위험, \quad ㄴ: 법률위험, \quad ㄷ: 계획위험
\end{choices}
\end{question}
\answer{4}
